\documentclass[11pt]{article}

% ============================================================================
% 1.1 - Analytic Stability Boundary for a Performative Market Maker
% LaTeX companion to 01-analytic-stability-boundary.md (REFLEX / math-theory).
% Self-contained and compilable with pdflatex.
% ============================================================================

\usepackage[margin=1in]{geometry}
\usepackage{amsmath,amssymb,amsthm}
\usepackage{booktabs}
\usepackage{xcolor}
\usepackage[colorlinks=true,linkcolor=blue,urlcolor=blue]{hyperref}
\usepackage{enumitem}

\newtheorem{theorem}{Theorem}
\newtheorem{corollary}{Corollary}
\newtheorem{assumption}{Assumption}

% Convenience macros matching the project's notation.
\newcommand{\eps}{\varepsilon}
\newcommand{\gbar}{\bar g}
\newcommand{\hdep}{h_{\mathrm{dep}}}
\newcommand{\href}{h_{\mathrm{ref}}}
\newcommand{\code}[1]{\texttt{#1}}
\newcommand{\BR}{\mathrm{BR}}

\title{\textbf{1.1 - Analytic Stability Boundary\\ for a Performative Market Maker}}
\author{REFLEX project - \code{math-theory}}
\date{}

\begin{document}
\maketitle

\noindent\textbf{STATUS: DONE (Priority 1).} This document discharges
methodology target \textbf{1.1}. It derives, in closed form, the
performative-prediction contraction modulus $m$ of the REFLEX
policy--distribution loop directly from the microstructure primitives of the
\code{endo\_market\_v2} simulator, and obtains the stability boundary
%
\begin{equation*}
  \text{stable} \quad\Longleftrightarrow\quad
  m = \frac{\eps\,\beta}{\gamma} < 1
  \quad\Longleftrightarrow\quad
  \eps < \frac{\gamma}{\beta},
\end{equation*}
%
with $\gamma$, $\beta$, and $\eps$ each given as a closed-form function of the
model parameters - none swept, none tuned. The result is then turned into a
falsifiable cross-check against the existing model-free estimator in
\code{analysis/response\_modulus.py}.

\paragraph{Why this is the headline contribution.}
Perdomo et al.\ (ICML 2020) prove that the repeated-retraining map contracts iff
$\eps < \gamma/\beta$, but treat $\gamma$ (strong convexity), $\beta$ (joint
smoothness) and $\eps$ (distribution sensitivity) as abstract Lipschitz
constants of an unspecified loss. For a structural market-making model these
constants are not free: they are pinned by the fill-intensity curve, the
adverse-selection channel, and the quoting-cost friction. We compute them. The
deliverable is an \emph{a-priori} stability boundary: predict whether a given
bond/dealer configuration converges before running a single RRM loop, then
verify the prediction against the model-free best-response (BR) slope probe.

% ===========================================================================
\section{Notation}

\begin{center}
\begin{tabular}{@{}lll@{}}
\toprule
Symbol & Reading & Definition / source \\
\midrule
$h$       & central half-spread (the policy coordinate) & bias $b_h$ of \code{LinearPolicy} \\
$h^*$     & self-consistent fixed-point half-spread     & root of the FOC, \S4 \\
$\hdep$   & deployed half-spread setting the toxic env. & \S1, A3 \\
$J(h;T)$  & expected one-step dealer objective          & \S2 \\
$\tau(h)$ & expected informed (toxic) notional          & \S2 \\
$T$       & frozen toxic level in one deployment, $T=\tau(\hdep)$ & \S1, A3 \\
$\eps$    & distribution sensitivity, $=|d\tau/dh|$     & \S3.1 \\
$\beta$   & joint smoothness, $=|\partial^2 J/\partial h\,\partial T|$ & \S3.2 \\
$\gamma$  & strong convexity, $=-\partial^2 J/\partial h^2$ & \S3.3 \\
$m$       & contraction modulus $=|\BR'(h^*)|=\eps\beta/\gamma$ & \S3.4 \\
$\rho$    & liquidity ratio \code{liquidity/liq\_mean}  & A2 \\
$g,\,u$   & mispricing $v-m$, standardized $u=g/\sigma_s$ & \S2 \\
$\gbar(u)$& gate mean (a 1-D Gaussian integral)         & \S2 \\
\bottomrule
\end{tabular}
\end{center}

% ===========================================================================
\section{Setup and assumptions}

We analyse the loop at the level of the \textbf{central half-spread} $h$ - the
bias coordinate $b_h$ of the dealer policy, which the convergence study
identifies as the dominant coordinate of the policy-to-distribution iterate map.
The model-free probe \code{response\_modulus.py} deploys constant-$h$, zero-skew
policies; we adopt the same restriction so the analytic and empirical moduli
measure the same object.

\begin{assumption}[single representative bond, zero skew]
One bond, $\mathrm{skew}=0$. The cross-sectional factor structure rescales but
does not change the scalar boundary; the multi-bond lift is Priority 1.5.
\end{assumption}

\begin{assumption}[reference state]
Curvature is evaluated at the probe's reference state $s_0=\code{simulator.reset}$
with inventory $q_0\approx 0$, mispricing $g=v-m$, and liquidity ratio
$\rho=\code{liquidity}/\code{liq\_mean}$. Expectations below are over the
client-flow noise at that state.
\end{assumption}

\begin{assumption}[frozen-environment best response]
This is the structural content of ``the operator is blind to $dD/d\phi$''. Within
a single deployment the learned operator $T_\theta$ conditions on a frozen
policy summary of the deployed regime. So when the dealer computes its best
response it sees the toxic-flow environment as fixed at the level induced by the
deployed spread $\hdep$ (write $T=\tau(\hdep)$), while the benign
demand-elasticity channel still responds to its candidate $h$. The performative
feedback enters only across deployments - exactly the cobweb the loop iterates.
\end{assumption}

\begin{assumption}[saturation non-binding]
The informed-flow cap $\code{info\_cap}\cdot\tanh(\cdot/\code{info\_cap})$ is in
its linear regime at the operating point. We carry the uncapped form and flag
the saturated regime as a documented corollary (\S5.3).
\end{assumption}

% ===========================================================================
\section{The expected single-step objective \texorpdfstring{$J(h)$}{J(h)}}

\subsection{Components from the simulator}
At zero skew, with $S$ = dealer sells (clients lift the ask), $B$ = dealer buys
(clients hit the bid), $q_{\mathrm{after}}=q_0+B-S$, $v$ = fundamental,
$v'$ = next fundamental, $m$ = mid:
%
\begin{align*}
  \text{spread\_capture}         &= (S+B)\,h, \\
  \text{inventory\_pnl}          &= q_{\mathrm{after}}\,(v'-v), \\
  \text{adverse\_selection\_loss}&= (B-S)\,(m-v).
\end{align*}
%
The dealer objective is
%
\begin{equation*}
  J = P\,\big[\, \text{spread\_capture} + \text{inventory\_pnl}
      - \text{adverse\_selection\_loss}
      - \code{inv\_risk\_weight}\cdot q_{\mathrm{after}}^2
      - w\,(h-\href)^2 \,\big],
\end{equation*}
%
with $w=\code{quote\_anchor\_weight}$, $\href=\code{quote\_anchor\_ref}$, and
$P:=\code{pnl\_scale}$.

\subsection{Uninformed (benign) flow}
Benign notional has the Gu\'eant--Lehalle--Fern\'andez-Tapia / Avellaneda--Stoikov
exponential fill intensity
%
\begin{equation*}
  U(h) = A\,e^{-kh}\,\rho,\qquad A=\code{base\_arrival\_rate},\;\; k=\code{demand\_elasticity},
\end{equation*}
%
split roughly $50/50$ into buys/sells with a zero-mean idiosyncratic imbalance,
so $\mathbb{E}[S_{\mathrm{uninf}}]=\mathbb{E}[B_{\mathrm{uninf}}]=\tfrac12 U(h)$
and $\mathbb{E}[S_{\mathrm{uninf}}-B_{\mathrm{uninf}}]=0$.

\subsection{Informed (toxic) flow}
With $\mathrm{signal}=g+\sigma_s\eta$, $\eta\sim N(0,1)$,
$\sigma_s=\code{info\_signal\_noise}$, the gate is
$\mathrm{gate}=\tanh(|\mathrm{signal}|/\sigma_s)$. The informed notional is
%
\begin{equation*}
  \tau(h) = \rho\,\gbar\,\big(I_b + \alpha f I\,e^{-c_t h}\big),
\end{equation*}
%
with $I_b=\code{info\_base\_intensity}$, $I=\code{info\_intensity}$,
$\alpha=\code{alpha}$, $f=\code{toxicity\_feedback}$,
$c_t=\code{info\_spread\_decay}$, and
$\gbar=\mathbb{E}_\eta[\tanh(|g+\sigma_s\eta|/\sigma_s)]$ the gate mean. Informed
traders pick the dealer off, so $\mathbb{E}[(B-S)_{\mathrm{inf}}]=-\tau(h)\,\mathrm{sign}(g)$
and the expected adverse-selection loss is
%
\begin{equation*}
  \mathbb{E}[\text{adverse}] = \psi\,\tau(h),\qquad \psi>0.
\end{equation*}

\subsection{Inventory terms vanish or are second order}
The fundamental increment $v'-v=\sigma_v\,\mathrm{shock}$ is drawn after, and
independently of, $q_{\mathrm{after}}$, so
$\mathbb{E}[\text{inventory\_pnl}]=\mathbb{E}[q_{\mathrm{after}}]\cdot 0=0$ for
every $h$. The quadratic risk $\mathbb{E}[q_{\mathrm{after}}^2]$ is $O(q_0^2)$
plus a flow-variance term that is second order at $q_0\approx 0$; we carry it as
a small additive curvature $\lambda_q\ge 0$ that only raises $\gamma$ (it is
stabilising) and never changes the boundary's form.

\subsection{The two Gaussian integrals}
Standardize $u=g/\sigma_s$. Because $\tanh$ is odd,
$\tanh(|x|)\,\mathrm{sign}(x)=\tanh(x)$, so the two state constants reduce to
clean 1-D Gaussian integrals in $u$ alone:
%
\begin{align*}
  \gbar(u) &= \mathbb{E}_{\eta\sim N(0,1)}\!\big[\tanh(|u+\eta|)\big], \\
  a(u)     &= \mathbb{E}_{\eta\sim N(0,1)}\!\big[\tanh(u+\eta)\big], \\
  \psi     &= \sigma_s\,u\,\frac{a(u)}{\gbar(u)} \;>\;0.
\end{align*}
%
Both are smooth, bounded, monotone-in-$|u|$ integrals evaluated by quadrature.
They are state constants (independent of $h$), so they pass through every
$h$-derivative untouched.

\subsection{Assembled objective}
Writing the toxic level as the frozen $T=\tau(\hdep)$ (A3):
%
\begin{equation}\label{eq:J}
  J(h;T) = P\,\Big[\,
    \underbrace{h\,A\,e^{-kh}\rho}_{\text{uninf.\ spread}}
    + \underbrace{h\,T}_{\text{toxic spread}}
    - \underbrace{\psi\,T}_{\text{adverse sel.}}
    - \underbrace{w\,(h-\href)^2}_{\text{quoting cost}}
    - \lambda_q \,\Big].
\end{equation}
%
\textbf{Key structural fact.} The adverse term $-\psi T$ is constant in $h$ (the
toxic level is frozen within a deployment), so it drops out of every
$h$-derivative below. It sets the level of profit - the echo-chamber gap of
Priority 1.2 - but \emph{adverse selection does not enter the stability
boundary}, only the optimal-vs-stable profit gap.

% ===========================================================================
\section{The best-response map and its slope}

\subsection{The map}
The dealer best-responds by maximising $J(h;T)$ at the frozen toxic level
$T=\tau(\hdep)$: $\BR(\hdep)=\arg\max_h J(h;\tau(\hdep))$. The RRM iteration is
$h_{k+1}=\BR(h_k)$; linearising at a fixed point $h^*$,
$h_{k+1}-h^*\approx \BR'(h^*)(h_k-h^*)$, so the loop contracts iff the modulus
$m:=|\BR'(h^*)|<1$.

\subsection{First-order condition and the three derivatives}
Define $F(h,T):=\partial J/\partial h$:
%
\begin{equation}\label{eq:FOC}
  F(h,T) = P\,\big[\, A\rho\,e^{-kh}(1-kh) + T - 2w(h-\href) \,\big],
  \qquad F=0 \text{ at the optimum.}
\end{equation}

\paragraph{(i) Curvature in $h$ (gives $\gamma$).}
Only the uninformed term and quoting cost depend on $h$ ($hT$ is linear):
%
\begin{align*}
  \frac{d}{dh}\!\big[A\rho e^{-kh}(1-kh)\big]
    &= A\rho e^{-kh}\big[-k(1-kh)-k\big]
     = A\rho e^{-kh}\,(k^2 h - 2k) \\
    &= A\rho k\,e^{-kh}(kh-2), \\
  \frac{d}{dh}\!\big[-2w(h-\href)\big] &= -2w,
\end{align*}
%
hence $\partial^2 J/\partial h^2 = P\,[\,A\rho k\,e^{-kh}(kh-2)-2w\,]$ and
%
\begin{equation}\label{eq:gamma}
  \boxed{\;
  \gamma := -\frac{\partial^2 J}{\partial h^2}
  = P\,\big[\, 2w + A\rho k\,e^{-kh}(2-kh) + \lambda_q \,\big].
  \;}
\end{equation}

\paragraph{(ii) Cross sensitivity (gives $\beta$).}
Only $hT$ carries a $T$-dependence in the gradient (coefficient $+1$); the
constant adverse term contributes nothing to $\partial J/\partial h$:
%
\begin{equation}\label{eq:beta}
  \boxed{\;
  \beta := \Big|\frac{\partial^2 J}{\partial h\,\partial T}\Big| = P = \code{pnl\_scale}.
  \;}
\end{equation}

\paragraph{(iii) Distribution sensitivity (gives $\eps$).}
\begin{equation}\label{eq:eps}
  \boxed{\;
  \eps := \Big|\frac{d\tau}{dh}\Big|
  = \rho\,\gbar\,\alpha f I\,c_t\,e^{-c_t h}.
  \;}
\end{equation}
$\eps$ is the slope of the toxic-flow response to the deployed spread (tighter
quotes summon more toxic flow) and is linear in $\alpha$ and in
$\code{toxicity\_feedback}=f$.

\subsection{Implicit-function theorem: the BR slope}
Differentiate $F(h^*(\hdep),\tau(\hdep))=0$ totally in $\hdep$, with
$F_h=-\gamma$, $F_T=+\beta$, $d\tau/d\hdep=-\eps$:
%
\begin{equation*}
  \frac{dh^*}{d\hdep}
  = -\frac{F_T\,(d\tau/d\hdep)}{F_h}
  = -\frac{(\beta)(-\eps)}{(-\gamma)}
  = -\frac{\beta\,\eps}{\gamma}.
\end{equation*}
%
So $\BR'(\hdep)=-\beta\eps/\gamma$ - a \emph{negative} slope. The cobweb is
oscillatory: deploying wider invites a narrower best response and vice versa.

\subsection{The modulus and the boundary}
\begin{theorem}[analytic stability boundary]\label{thm:boundary}
At the fixed point $h^*$,
%
\begin{equation}\label{eq:modulus}
  m(h^*) = |\BR'(h^*)| = \frac{\eps\,\beta}{\gamma}
  = \frac{\rho\,\gbar\,\alpha f I\,c_t\,e^{-c_t h^*}}
         {2w + A\rho k\,e^{-k h^*}(2-k h^*) + \lambda_q},
\end{equation}
%
(the common $P=\code{pnl\_scale}$ cancels), and the loop is stable iff
%
\begin{equation*}
  \eps < \frac{\gamma}{\beta}
  \;\Longleftrightarrow\;
  \rho\gbar\alpha f I c_t e^{-c_t h^*}
  < 2w + A\rho k\,e^{-k h^*}(2-k h^*) + \lambda_q.
\end{equation*}
\end{theorem}
%
Every quantity is a closed form in the config and $h^*$, so the inequality is an
a-priori predictor. $m=|\BR'|$ is exactly what \code{response\_modulus.py} returns,
so \eqref{eq:modulus} is the directly comparable prediction (\S7).

% ===========================================================================
\section{Locating the fixed point \texorpdfstring{$h^*$}{h*}}

The boundary is evaluated where the deployed spread is its own best response.
Setting $F(h,\tau(h))=0$ with the toxic level evaluated self-consistently:
%
\begin{equation}\label{eq:fixedpoint}
  A\rho\,e^{-k h^*}(1-k h^*) + \tau(h^*) - 2w(h^*-\href) = 0,
  \qquad \tau(h^*)=\rho\gbar\big(I_b+\alpha f I e^{-c_t h^*}\big).
\end{equation}
%
A smooth scalar equation; solve by Newton or bisection on
$[0,\code{max\_half\_spread}]$. A unique interior root exists wherever
$\gamma(h)>0$ (strict concavity of $J$ in $h$) - exactly the stability
condition holding with room to spare.

% ===========================================================================
\section{Corollaries (each a falsifiable prediction)}

\begin{corollary}[linear-in-adversariality scaling]
At fixed $h^*$, $m$ is linear in $\alpha f$, and $m=1$ at
\begin{equation*}
  (\alpha f)^* = \frac{\gamma}{\rho\,\gbar\,I\,c_t\,e^{-c_t h^*}},
\end{equation*}
predicting the \code{toxicity\_feedback} value at which \code{sweep\_feedback.yaml}
crosses the boundary before the sweep is run.
\end{corollary}

\begin{corollary}[lazy-deploy raises effective curvature]
Under \code{update\_rule="rgd"} with step $\eta$, one gradient step gives the map
$M(h)=h+\eta\,G(h)$ with $G(h)=\partial_h J(h;\tau(h))$ and linearised modulus
$m_{\mathrm{RGD}}=|1-\eta(\gamma-\eps\beta)|$ in the worst-case (destabilising)
orientation. Contraction requires $\gamma_{\mathrm{eff}}:=\gamma-\eps\beta>0$,
i.e.\ the same boundary $\eps<\gamma/\beta$. Taking $K$ lazy steps compounds the
contraction, enlarging the stable region.
\end{corollary}

\begin{corollary}[the modulus saturates rather than blows up]
As $h^*$ grows, $A\rho k e^{-k h^*}(2-k h^*)\to 0$, so $\gamma\to 2wP$ (a floor,
not zero) while $\eps$ also decays. Their ratio approaches a finite plateau
$m_\infty\approx \eps(h^*)/(2w)$ rather than diverging - the analytic
explanation of the empirically observed ``modulus saturates ($\sim 1.25$) past
the boundary''. The GLFT term even turns negative for $k h^*>2$ (defensive
widening, low curvature).
\end{corollary}

\begin{corollary}[why $\eps$, not $\alpha$, is the clean control variable]
Sweeping $\alpha$ moves both $\eps$ (numerator, $\propto\alpha$) and $h^*$ (the
dealer widens), feeding back through $e^{-c_t h^*}$ and the $\gamma$ curvature
term, so $m(\alpha)$ is confounded and its sign is not robust to universe size.
Sweeping $f=\code{toxicity\_feedback}$ at fixed structural regime moves $\eps$
linearly with $h^*$ far less perturbed, isolating the performative channel.
\end{corollary}

% ===========================================================================
\section{Symbol \texorpdfstring{$\to$}{->} config map}

\begin{center}
\begin{tabular}{@{}llll@{}}
\toprule
Symbol & Meaning & Config field & Default \\
\midrule
$A$        & uninformed arrival scale          & \code{clients.base\_arrival\_rate} & 1.0 \\
$k$        & demand elasticity (GLFT decay)    & \code{clients.demand\_elasticity}  & 1.5 \\
$I_b$      & toxic baseline level              & \code{clients.info\_base\_intensity} & 0.60 \\
$I$        & toxic slope scale                 & \code{clients.info\_intensity}     & 1.4 \\
$\alpha$   & adversariality                    & \code{clients.alpha}               & 0.15 \\
$f$        & performative feedback gain        & \code{clients.toxicity\_feedback}  & 0.22 \\
$c_t$      & toxic spread-decay                & \code{clients.info\_spread\_decay} & 1.5 \\
$\sigma_s$ & informed signal noise             & \code{clients.info\_signal\_noise} & 0.6 \\
$w$        & quoting-cost convexity            & \code{reward.quote\_anchor\_weight}& 0.25 \\
$\href$    & quoting-cost anchor               & \code{reward.quote\_anchor\_ref}   & 1.0 \\
$P$        & global P\&L scale ($=\beta$)       & \code{reward.pnl\_scale}           & 1.0 \\
$\lambda_q$& inventory-risk curvature          & \code{reward.inv\_risk\_weight}    & 0.05 \\
$\rho$     & liquidity ratio at reference      & state (\code{liquidity/liq\_mean}) & $\approx 1$ \\
$\gbar,\psi$ & gate mean / adverse severity    & derived (Gaussian integrals, \S3.5)& - \\
\bottomrule
\end{tabular}
\end{center}

% ===========================================================================
\section{Validation protocol (predict then verify)}

\begin{enumerate}[leftmargin=*]
\item \textbf{Predict.} For each $f$ in \code{sweep\_feedback.yaml}, solve
  \eqref{eq:fixedpoint} for $h^*$, then evaluate $m_{\mathrm{pred}}(f)=\eps\beta/\gamma$
  from \eqref{eq:gamma}, \eqref{eq:beta}, \eqref{eq:eps}.
\item \textbf{Measure.} Run \code{measure\_response\_modulus} at $\href\approx h^*$
  for the common-random-number BR-slope estimate $m_{\mathrm{meas}}(f)$.
\item \textbf{Compare.} Report $m_{\mathrm{pred}}$ vs $m_{\mathrm{meas}}$ across the
  sweep with cross-seed IQR bands. The claim: the predicted boundary $(\alpha f)^*$
  coincides with the measured crossing $m_{\mathrm{meas}}=1$, and
  $m_{\mathrm{pred}}\approx m_{\mathrm{meas}}$ in the contracting regime, diverging
  only in the saturated regime (Corollary 3).
\end{enumerate}
%
A reference implementation belongs in
\code{endo\_market\_v2/endo\_market/analysis/analytic\_boundary.py} (a pure
function of \code{Config}), reusing \code{response\_modulus.py} for step 2.

% ===========================================================================
\section{Honest caveats}

\begin{itemize}[leftmargin=*]
\item \textbf{$\beta=\code{pnl\_scale}$ is exact only under the frozen-environment
  FOC (A3).} A future operator modelling $dD/d\phi$ (PerfGD, Priority 1.2) makes
  $h\tau(h)$ $h$-responsive and $\beta$ gains a $-h c_t$ correction; the boundary
  form survives but $\beta$ is no longer a bare constant. This is the bridge to
  1.2, not a defect.
\item \textbf{$\lambda_q,\gbar,\psi$ are state-dependent.} The boundary is stated
  at the probe reference state; the phase diagram (Priority 1.5) reports it as a
  band over the reference-state distribution.
\item \textbf{The cap (A4) is assumed slack.} Deep in the unstable regime the cap
  re-convexifies the problem; the analytic modulus is valid up to and slightly
  past the boundary, not arbitrarily far into divergence.
\item \textbf{Single bond, zero skew (A1).} The scalar boundary is the base case;
  the factor-model lift (1.5) and the multi-dealer $\eps<\gamma/(N\beta)$ (1.3)
  build on it but are out of scope here.
\end{itemize}

% ===========================================================================
\section*{References}
\begin{itemize}[leftmargin=*]
\item J.\ Perdomo, T.\ Zrnic, C.\ Mendler-D\"unner, M.\ Hardt.
  \emph{Performative Prediction.} ICML 2020. - the $m=\eps\beta/\gamma$
  criterion this work makes structural.
\item M.\ Avellaneda, S.\ Stoikov. \emph{High-frequency trading in a limit-order
  book.} Quant.\ Finance 2008; O.\ Gu\'eant, C.-A.\ Lehalle, J.\
  Fern\'andez-Tapia (GLFT), 2013. - the exponential fill intensity
  $\lambda(\delta)=A e^{-k\delta}$ underlying $\gamma$.
\item See \code{../references.bib} and \code{literature/literature-raghav/README.md}
  for the full citation map (Barzykin et al.\ 2025 for the $d\tau/dh$ grounding;
  Cont--Kukanov--Stoikov for the informed-flow slope used in the triangulation
  extension of 1.1).
\end{itemize}

\end{document}
